% !TeX root = ../../Book5.tex
% 纲要标题：### B.2 簇代数的变异与小秩例子
% 纲要位置：工作记录/计划纲要.md，第 4585 行。
% #### B.2.1 种子与交换关系
% #### B.2.2 二阶五周期的完整计算
% #### B.2.3 Laurent 现象与计算范围
\section{簇代数的变异与小秩例子}
\label{b5:sec:B02}

簇代数通过一串有理变量替换生成交换环。每次替换只改变一个变量，新变量的分子是两个单项式的和。这些替换不断引入分母，但在某些小秩例子中会出现有限循环，变量仍能写成 Laurent 多项式。下面从一个二阶例子说明这一现象。

\subsection{种子与交换关系}
\label{b5:subsec:B02:01}

先取特征零域 \(k\)、有理函数域 \(k(x_1,\ldots,x_n)\)，以及反对称整数矩阵 \(B=(b_{ij})\)。向量 \(\mathbf x=(x_1,\ldots,x_n)\) 的分量代数独立，这一对数据称为一个无系数种子。沿位置 \(r\) 的变异保留 \(x_i\)、\(i\ne r\)，并用
\[
x_rx'_r=\prod_{b_{ir}>0}x_i^{b_{ir}}
+\prod_{b_{ir}<0}x_i^{-b_{ir}}
\]
定义 \(x'_r\)。空乘积取一。矩阵改为
\[
b'_{ij}=
\begin{cases}
-b_{ij},&i=r\text{ 或 }j=r,\\
b_{ij}+[b_{ir}]_+[b_{rj}]_+-[-b_{ir}]_+[-b_{rj}]_+,&i,j\ne r,
\end{cases}
\quad [a]_+=\max(a,0).
\]
全部有限次变异得到的变量称为簇变量，它们在原有理函数域中生成的 \(k\) 子代数称为此种子的簇代数。
\ProseParagraphBreak
两次在相同位置变异会回到原种子。因为第 \(r\) 列变号，两交换单项式互换，分子不变，故 \(x''_r=x_r\)。矩阵的两项修正也互换符号，相消后恢复原值；同时 \(b'_{ji}=-b'_{ij}\)，所以反对称性保持。这个可逆有理替换还说明新变量仍然代数独立并生成同一有理函数域。
这里采用无冻结变量、反对称矩阵的基本版本；带系数、冻结变量和反对称可对称化矩阵的其他版本还需补充相应定义。

\subsection{二阶五周期的完整计算}
\label{b5:subsec:B02:02}

取 \(B=\begin{psmallmatrix}0&1\\-1&0\end{psmallmatrix}\)，沿两个位置交替变异。令连续簇变量为 \(z_1=x,z_2=y\)，则
\(z_{i-1}z_{i+1}=1+z_i\)。
直接计算得
\[
z_3=\frac{1+y}{x},\qquad
z_4=\frac{1+x+y}{xy},\qquad
z_5=\frac{1+x}{y},\qquad
z_6=x,\quad z_7=y.
\]
例如
\(1+z_4=(x+1)(y+1)/(xy)\)，除以 \(z_3\) 得到 \(z_5\)；再继续两步就回到原来的两变量。任意变异串可消去相邻两次相同位置的变异；两个位置之间剩下的串必然交替，所以这个递推穷尽了所有可能。五个所列有理式两两不同，因此恰有五个可能簇变量，并有五个无序簇 \(\{z_i,z_{i+1}\}\)，下标模五。这里五周期指无序簇；保留位置编号的种子还须考虑两个位置的交换。它们都为原变量的 Laurent 多项式，但一般不是普通多项式。
\ProseParagraphBreak
原始论文中的二阶递推见\cite[Example 2.5; Section 6]{FominZelevinsky2002ClusterFoundations}，本节选择了其中系数一、两个指数都为一的无系数 \(A_2\) 情形。

\subsection{Laurent 现象与计算范围}
\label{b5:subsec:B02:03}

\begin{theorem}[Laurent 现象]\label{b5:thm:B-laurent}
设 \(k\) 为特征零域，\(n\ge1\) 为整数，\(x_1,\ldots,x_n\) 在 \(k\) 上代数独立，\(\mathcal F=k(x_1,\ldots,x_n)\)，\(B\in M_n(\mathbb Z)\) 为反对称交换矩阵。以 \(\Sigma=((x_1,\ldots,x_n),B)\) 为 \(\mathcal F\) 中没有系数、没有冻结变量的初始种子，按无系数种子变异产生的每个簇变量都属于
\(k[x_1^{\pm1},\ldots,x_n^{\pm1}]\)。
\end{theorem}
\begin{proofsketch}
证明不能只对路径长度直接使用交换式，因为除以一个已知 Laurent 多项式可能带来非单项式分母。解决这一点的办法，是把有限变异路径与每个路径顶点的邻边一起放进一棵有限树，并同时比较几条较短路径所得的表达。

先把交换式的系数提升到通用系数的多项式环：初始交换多项式及其变换关系所需的独立系数暂作不定元，最后统一代入一。对连续标号 \(i,j,i\) 的三步变异，省略未改变的变量，写成
\[
 z=\frac{P(y)}x,\qquad u=\frac{Q(z)}y,\qquad v=\frac{R(u)}z.
\]
矩阵变异规定的指数变换给出交换多项式的相容关系
\(R(U)=C(U)P\bigl(Q(0)/U\bigr)\)，其中 \(C\) 对 \(U\) 是多项式，对未改变变量是 Laurent 多项式。于是
\[
 v=\frac{R\bigl(Q(z)/y\bigr)-R\bigl(Q(0)/y\bigr)}z
     +xC\bigl(Q(0)/y\bigr).
\]
第一项的分子是关于 \(z\) 的多项式且常数项为零，所以能够除以 \(z\) 而不增加分母。这先证明三步后的变量仍在初始 Laurent 环中。

在通用系数环中还可证明 \(\gcd(z,u)=\gcd(z,v)=1\)。前一式由
\(u\equiv Q(0)/y\pmod z\) 得到，\(P\) 与 \(Q(0)\) 的独立系数使公因子只能是单位。后一式将上述差商模 \(z\) 化成
\(R'\bigl(Q(0)/y\bigr)Q'(0)/y+xC\bigl(Q(0)/y\bigr)\)，再利用 \(C\) 与 \(P\) 的独立系数。现在对有限树的路径长度归纳：某个目标变量沿两条较短路径，分别写成分母为 \(z^a\) 和 \(u^bv^c\) 的分式，分子均在初始 Laurent 环中。这个环是唯一分解环，而两分母互素，故该变量没有非单位分母。

这里省略了通用系数的统一选取、相容关系的逐指数核验，以及有限树归纳的顶点记号；它们分别保证系数特殊化合法、三步消去可在每个位置实行、任意有限路径都被覆盖。完整论证见\cite[Section 3, Theorem 3.2 and Lemma 3.3]{FominZelevinsky2002ClusterFoundations}。归纳所得系数是通用参数的多项式，因此参数代入一后仍只有初始变量的单项式分母，给出本节的无系数版本。
\end{proofsketch}

即使交换时分子越来越复杂，约掉因子后的分母仍只有单项式；前一小节的 \(A_2\) 计算是这一结论的具体实例。
\ProseParagraphBreak
Laurent 性也不自动意味着只有有限多个簇变量。例如改成二阶递推
\(z_{i-1}z_{i+1}=1+z_i^2\)，会出现不同的行为。有限型分类需要控制全部变异，不是看到前几步没有新变量就能断定。类似地，正系数性、规范基、量子化和表示论实现分别是额外问题，不能都从 Laurent 现象本身推出。
